\uuid{qTo2}
\exo7id{4227}
\titre{exo7 4227}
\auteur{quercia}
\organisation{exo7}
\datecreate{2010-03-12}
\isIndication{false}
\isCorrection{false}
\chapitre{Calcul d'intégrales}
\sousChapitre{Théorie}
\cpge{mpsi-mp2i.raisonnement-vocabulaire-ensembliste.c, pcsi.raisonnement-vocabulaire-ensembliste.d, ptsi.raisonnement-vocabulaire-ensembliste.d, mpsi-mp2i.fonctions-valeurs-reelles-complexes.b, pcsi.techniques-calcul-fonctions.b, ptsi.techniques-calcul-fonctions.b, mpsi-mp2i.integration.c, pcsi.integration.b, ptsi.integration.b, bcpst1.derivees-primitives-integrales.d, bcpst1.limites-continuite.c, bcpst1.limites-continuite.d, bcpst1.integration-segment.a, bcpst1.integration-segment.d}
\module{Analyse}
\niveau{L1}
\difficulte{}

\contenu{
\question{
Soit $f : {[a,b]} \to {[c,d]}$ continue, bijective, strictement croissante.

Calculer $ \int_{t=a}^b f(t)\,d t +  \int _{u=c}^d f^{-1}(u)\,d u$
(faire un dessin, et commencer par le cas où $f$ est de classe $\mathcal{C}^1$).
}
}
